\section{Erd\H{o}s Problem 41 (Round 4)}

\subsection{Problem}
An infinite set $A\subset \mathbb{N}$ is a \emph{$B_3$-set} if all sums
\[
a_1+a_2+a_3\qquad (a_1\le a_2\le a_3,\ a_i\in A)
\]
are distinct. The target statement is
\[
\liminf_{N\to\infty}\frac{|A\cap[1,N]|}{N^{1/3}}=0.
\]
Write
\[
A(N):=|A\cap[1,N]|.
\]

\subsection{1) ROUND-4 OBJECTIVE}
\textbf{Path (C): obstruction/correction.}

Round~3 showed that the Round~2 ``full proof'' was not gap-free, and reduced the original Erd\H{o}s problem to the oscillatory regime
\[
0<\liminf_{N\to\infty}\frac{A(N)}{N^{1/3}}
<
\limsup_{N\to\infty}\frac{A(N)}{N^{1/3}}
<\infty.
\]
The most promising Round~4 advance is therefore not to restart the whole proof, but to sharpen that obstruction: we prove that no such oscillation can stabilize on an asymptotically dense multiplicative scale. This rules out an entire class of possible counterexamples.

\subsection{2) Round-3 FOUNDATION USED}
We use the following vetted results from Round~3.
\begin{itemize}
  \item \textbf{Round~1 Lemma 41.1 (retained in Round~3).} For every infinite $B_3$-set,
  \[
  A(N)\le (18N)^{1/3}+2,
  \]
  hence
  \[
  \limsup_{N\to\infty}\frac{A(N)}{N^{1/3}}<\infty.
  \]
  \item \textbf{Round~3 Proposition 5.1.} If the limit
  \[
  \lim_{N\to\infty}\frac{A(N)}{N^{1/3}}
  \]
  exists for an infinite $B_3$-set, then it must equal $0$. Equivalently, no infinite $B_3$-set satisfies
  \[
  A(N)\sim \alpha N^{1/3}
  \qquad (\alpha>0).
  \]
  \item \textbf{Round~3 Proposition 5.2.} Given only the verified inputs, the original Erd\H{o}s problem is equivalent to ruling out the oscillatory regime
  \[
  0<\liminf_{N\to\infty}\frac{A(N)}{N^{1/3}}
  <
  \limsup_{N\to\infty}\frac{A(N)}{N^{1/3}}
  <\infty.
  \]
\end{itemize}

\subsection{3) NEW INSIGHT / TOOL (ROUND-4)}
\textbf{Dense-subsequence rigidity by monotonicity.}

Let $B(N)$ be any nondecreasing counting function. If $B(N)N^{-\gamma}$ converges along a sequence $N_k$ with $N_{k+1}/N_k\to 1$, then in fact $B(N)N^{-\gamma}$ converges along \emph{all} integers $N$. Applied with $\gamma=1/3$, this means that a $B_3$-set cannot even look pseudo-cubic along an asymptotically dense multiplicative subsequence.

This is genuinely new in Round~4: Round~3 ruled out only \emph{global} pseudo-cubic asymptotics; here we rule out \emph{subsequential} pseudo-cubic asymptotics on scales $N_{k+1}=(1+o(1))N_k$.

\subsection{4) ATTACK PLAN (ROUND-4)}
The exact gap left after Round~3 is the possibility of a $B_3$-set with
\[
0<\liminf f(N)<\limsup f(N)<\infty,
\qquad f(N):=\frac{A(N)}{N^{1/3}},
\]
but with no full limit.

Round~4 proves the following claims.
\begin{enumerate}
  \item A general rigidity lemma: if $f(N_k)$ converges along a sequence with $N_{k+1}/N_k\to 1$, then $f(N)$ converges globally.
  \item Using Round~3 Proposition~5.1, this implies that no infinite $B_3$-set admits a subsequence $N_k$ with $N_{k+1}/N_k\to 1$ and $f(N_k)\to \alpha>0$.
  \item Therefore any hypothetical counterexample must oscillate on every asymptotically dense multiplicative scale; equivalently, every positive cluster value of $f$ can only be attained along multiplicatively sparse subsequences.
\end{enumerate}
This does not settle the full problem, but it excludes a broad family of oscillatory profiles that remained possible after Round~3.

\subsection{5) WORK (ROUND-4)}

\paragraph{Lemma 5.1 (dense-subsequence rigidity).}
Let $B:\mathbb{N}\to[0,\infty)$ be nondecreasing, let $\gamma>0$, and define
\[
g(N):=\frac{B(N)}{N^{\gamma}}.
\]
Suppose there is an increasing sequence $N_1<N_2<\cdots$ such that
\[
\frac{N_{k+1}}{N_k}\to 1
\qquad\text{and}\qquad
g(N_k)\to \alpha
\]
for some $\alpha\in[0,\infty)$. Then
\[
g(N)\to \alpha\qquad (N\to\infty).
\]

\paragraph{Proof.}
Fix large $N$, and choose the unique index $k=k(N)$ such that
\[
N_k\le N < N_{k+1}.
\]
Since $B$ is nondecreasing,
\[
B(N_k)\le B(N)\le B(N_{k+1}).
\]
Dividing by $N^{\gamma}$ gives
\[
g(N_k)\Bigl(\frac{N_k}{N}\Bigr)^{\gamma}
\le g(N)
\le g(N_{k+1})\Bigl(\frac{N_{k+1}}{N}\Bigr)^{\gamma}.
\]
Because $N_k\le N < N_{k+1}$, we also have
\[
\frac{N_k}{N}\ge \frac{N_k}{N_{k+1}}
\qquad\text{and}\qquad
\frac{N_{k+1}}{N}\le \frac{N_{k+1}}{N_k}.
\]
Hence
\[
g(N_k)\Bigl(\frac{N_k}{N_{k+1}}\Bigr)^{\gamma}
\le g(N)
\le g(N_{k+1})\Bigl(\frac{N_{k+1}}{N_k}\Bigr)^{\gamma}.
\]
Now $N\to\infty$ implies $k(N)\to\infty$, and by hypothesis
\[
g(N_k)\to \alpha,
\qquad
g(N_{k+1})\to \alpha,
\qquad
\frac{N_{k+1}}{N_k}\to 1.
\]
Therefore both the lower and upper bounds tend to $\alpha$. By the squeeze theorem,
\[
g(N)\to \alpha.
\]
This proves the lemma.

\paragraph{Corollary 5.2 (no positive asymptotic on asymptotically dense scales).}
Let $A$ be an infinite $B_3$-set and define
\[
f(N):=\frac{A(N)}{N^{1/3}}.
\]
Then there is no increasing sequence $N_1<N_2<\cdots$ such that
\[
\frac{N_{k+1}}{N_k}\to 1
\qquad\text{and}\qquad
f(N_k)\to \alpha>0.
\]

\paragraph{Proof.}
Assume such a sequence exists. Since $A(N)$ is nondecreasing, Lemma~5.1 applies with $B=A$ and $\gamma=1/3$, and yields
\[
\frac{A(N)}{N^{1/3}}\to \alpha.
\]
But Round~3 Proposition~5.1 says that if the full limit $\lim_{N\to\infty}A(N)N^{-1/3}$ exists for an infinite $B_3$-set, then it must equal $0$. This contradicts $\alpha>0$. Therefore no such sequence can exist.

\paragraph{Corollary 5.3 (positive cluster values are multiplicatively sparse).}
Let $A$ be an infinite $B_3$-set and let $\alpha>0$. If $N_k$ is any increasing sequence satisfying
\[
\frac{A(N_k)}{N_k^{1/3}}\to \alpha,
\]
then
\[
\frac{N_{k+1}}{N_k}\not\to 1.
\]
Equivalently,
\[
\limsup_{k\to\infty}\frac{N_{k+1}}{N_k}>1.
\]

\paragraph{Proof.}
If $N_{k+1}/N_k\to 1$, then Corollary~5.2 gives a contradiction. So $N_{k+1}/N_k\not\to 1$. Since every ratio is at least $1$, this is equivalent to
\[
\limsup_{k\to\infty}\frac{N_{k+1}}{N_k}>1.
\]

\paragraph{Proposition 5.4 (sharpened obstruction to a counterexample).}
Assume, for contradiction, that the original Erd\H{o}s statement fails for some infinite $B_3$-set $A$. Then, by Round~3 Proposition~5.2,
\[
0<\ell:=\liminf_{N\to\infty}f(N)
<
L:=\limsup_{N\to\infty}f(N)
<\infty.
\]
Let
\[
K:=\Bigl\{\beta\in\mathbb{R}:\ \exists N_k\to\infty\text{ with }f(N_k)\to\beta\Bigr\}
\]
be the cluster set of $f$. Then $K\subset[\ell,L]\subset(0,\infty)$, and for every $\alpha\in K$, every increasing sequence $N_k$ with $f(N_k)\to\alpha$ must be multiplicatively sparse in the sense of Corollary~5.3.

In particular, no hypothetical counterexample can stabilize near any positive level $\alpha$ on a scale $N_{k+1}=(1+o(1))N_k$.

\paragraph{Proof.}
The inequality
\[
0<\ell<L<\infty
\]
is exactly the unresolved regime identified in Round~3 Proposition~5.2. By definition of cluster set, every $\alpha\in K$ is obtained as the limit of $f(N_k)$ along some sequence $N_k\to\infty$. Since $K\subset[\ell,L]\subset(0,\infty)$, Corollary~5.3 applies to every such sequence and shows that $N_{k+1}/N_k\not\to1$. This is precisely the asserted multiplicative sparseness, and the final sentence is just a restatement.

\subsection{6) ADVERSARIAL VERIFICATION}
We now test the new argument for weaknesses.
\begin{itemize}
  \item \textbf{Use of Round~3 results.} The only inherited nontrivial input beyond monotonicity is Round~3 Proposition~5.1, i.e. the verified ``no pseudo-cubes'' consequence of Helm. We do \emph{not} reuse the unsupported stronger oscillation theorem from Round~2.
  \item \textbf{Hidden regularity assumptions.} Lemma~5.1 uses only that $B(N)$ is nondecreasing. No convexity, differentiability, or regular variation is assumed.
  \item \textbf{Boundary case $\alpha=0$.} Lemma~5.1 remains true for $\alpha=0$. Corollary~5.2 only excludes $\alpha>0$, exactly because Round~3 Proposition~5.1 leaves $\alpha=0$ possible.
  \item \textbf{Could the ratios oscillate?} Yes, and our statement allows that. Corollary~5.3 says only that the ratios cannot converge to $1$ if a positive cluster value is present. This is the precise conclusion proved.
  \item \textbf{Does this settle the original Erd\H{o}s problem?} No. A counterexample could still oscillate between positive levels along sufficiently sparse subsequences. The remaining open regime is now narrower: any such oscillation must occur on genuinely macroscopic multiplicative jumps.
\end{itemize}

\subsection{7) FINAL}
\textbf{UNRESOLVED (BUT STRICTLY ADVANCED).}

Round~3 reduced the problem to the regime
\[
0<\liminf_{N\to\infty}\frac{A(N)}{N^{1/3}}
<
\limsup_{N\to\infty}\frac{A(N)}{N^{1/3}}
<\infty.
\]
Round~4 proves a new rigidity theorem showing that such a regime cannot stabilize along any asymptotically dense multiplicative subsequence $N_{k+1}=(1+o(1))N_k$. Equivalently, every positive cluster value of $A(N)/N^{1/3}$ must be realized only along multiplicatively sparse subsequences. This rules out an entire class of possible counterexamples and is strictly stronger than any result established in earlier rounds.

\subsection{8) COMPLETION ESTIMATE}
\textbf{COMPLETION: 78\%}

\subsection{9) REFERENCES}
\begin{itemize}
  \item P. Erd\H{o}s, \emph{Problems and results on combinatorial number theory III}, in \emph{Number Theory Day} (Rockefeller Univ., 1976), Lecture Notes in Math. \textbf{626}, Springer, 1977, pp.~43--72.
  \item M. Helm, \emph{On the Distribution of $B_3$-Sequences}, J. Number Theory \textbf{58} (1996), 124--129.
  \item K. O'Bryant, \emph{A Complete Annotated Bibliography of Work Related to Sidon Sequences}, Electron. J. Combin. \textbf{DS11} (2004). Used for the recorded abstract of Helm's paper.
  \item J. Cilleruelo, \emph{Conjuntos de Sidon}, in \emph{Aritmética, Grupos y Análisis}, AGRA II lecture notes (2015; later online circulation checked in 2022). Used for the statement that the odd-$h$ liminf problem remains open there, and that Helm proves only the pseudo-cube obstruction.
\end{itemize}
